\section{Гипотеза ван ден Эссена и явные координаты ядер LND}

\subsection{Точная гипотеза и её логическое место}

Работа A.~van den Essen, P.~van Rossum, \emph{Triangular derivations related
to problems on affine $n$-space}, Proc. Amer. Math. Soc. \textbf{130}
(2002), 1311--1322, DOI
\href{https://doi.org/10.1090/S0002-9939-01-06212-8}{10.1090/S0002-9939-01-06212-8},
содержит соответствующую гипотезу как Conjecture~5.1, с.~1319.
Точная формулировка использует
\[
 f(u)=u^3-3u,\qquad g(u)=u^4-4u^2,\qquad h(u)=u^5-10u,
\]
\[
 A=\mathbb C[t,u,x,y,z],\qquad
 D=t\partial_u+f'(u)\partial_x+g'(u)\partial_y+h'(u)\partial_z.
\]
Утверждается, что $\ker D\not\cong_{\mathbb C}\mathbb C^{[4]}$.
Именно эту конкретную гипотезу рассматриваем здесь под названием
\emph{гипотеза E}. Множитель $t$ перед $\partial_u$ существенен.

В первоисточнике Theorem~4.1 и Corollary~4.2, с.~1317--1318, связывают
выпрямляемость вложения с полиномиальностью ядра. Наличие слайса
обсуждается в Theorem~3.1. Поэтому проверенный выпрямляющий автоморфизм
трилистника влечёт отрицание именно Conjecture~5.1. Это не является
доказательством общей проблемы сокращения: кандидат имеет размерность
$4$, а после добавления слайса получается полиномиальное кольцо
размерности $5$.

\subsection{Элементарный конструктивный перенос}

Следующее доказательство выписано непосредственно, чтобы зависимость от
проверяемого выпрямителя была полностью прозрачной. Сам принцип переноса
не заявляется как новый: он согласуется с указанным Theorem~4.1.

Пусть $k$ --- поле характеристики нуль,
$\gamma(u)=(f_1(u),f_2(u),f_3(u))$, и дан полиномиальный автоморфизм
\[
 F=(P,Q,R),\qquad F(\gamma(u))=(u,0,0),\qquad G=F^{-1}.
\]
Пусть $0\ne b(t)\in k[t]$. Рассмотрим
\[
 D_b=b(t)\partial_u+\sum_{i=1}^3f_i'(u)\partial_{x_i}
 \quad\text{на }A=k[t,u,x_1,x_2,x_3].
\]
Обозначим $V_i=f_i(u)-b(t)x_i$, $V=(V_1,V_2,V_3)$ и положим
\begin{equation}\label{eq:lnd-coordinates}
 p=P(V),\qquad q=\frac{Q(V)}{b(t)},\qquad
 r=\frac{R(V)}{b(t)},\qquad s=\frac{u-P(V)}{b(t)}.
\end{equation}
Тогда все четыре выражения являются полиномами и
\begin{equation}\label{eq:lnd-conjugacy}
 A=k[t,p,q,r,s],\qquad D_b=\partial_s,\qquad
 \ker D_b=k[t,p,q,r].
\end{equation}

\paragraph{Доказательство полиномиальности.}
Для любого полинома $L$ разность $L(\gamma(u)-b(t)X)-L(\gamma(u))$
делится на $b(t)$. Применив это к $P,Q,R$ и использовав
$P(\gamma)=u$, $Q(\gamma)=R(\gamma)=0$, получаем точную делимость
числителей в \eqref{eq:lnd-coordinates}. Это рассуждение происходит
в кольце полиномов; простота корней $b$ не требуется.

\paragraph{Явная обратная замена.}
Считая $t,p,q,r,s$ независимыми переменными, определим
\begin{equation}\label{eq:lnd-inverse}
 u=p+b(t)s,\qquad
 x_i=\frac{f_i(p+b(t)s)-G_i(p,b(t)q,b(t)r)}{b(t)}.
\end{equation}
Из $G(p,0,0)=\gamma(p)$ следует, что числитель также делится на
$b(t)$. В частности, формулы действуют при кратных корнях $b$.
Подстановка \eqref{eq:lnd-inverse} даёт
$V=G(p,b(t)q,b(t)r)$; последующее применение $F$ восстанавливает
$(p,b(t)q,b(t)r)$. Обратная композиция столь же непосредственно
восстанавливает $u$ и $X$. Следовательно, это взаимно обратные
полиномиальные замены координат.

Наконец, $D_b(t)=D_b(V_i)=0$, поэтому $D_b(p)=D_b(q)=D_b(r)=0$,
а $D_b(s)=1$. Это доказывает \eqref{eq:lnd-conjugacy}.

\paragraph{Что именно даёт применение.}
Для каждой модели с удостоверенным автоморфизмом $F$ получаем не только
абстрактную полиномиальность ядра, но явное сопряжение бесконечного
семейства $D_b$ с переносом. Ранг каждого такого ненулевого LND равен
$1$: в одной системе координат он уничтожает четыре координаты и
дифференцирует пятую. Новое применение к конкретной модели следует
признавать только после проверки обоих полиномиальных автоморфизмов
$F$ и $G$; проверка одного значения $F(\gamma(u))$ недостаточна.

\subsection{Слайс существует и без выпрямления}

Для трилистника имеется простой полином восстановления параметра
\[
 F_0(x,y,z)=yz-x^3-5xy+2z-7x,\qquad
 F_0(f(u),g(u),h(u))=u.
\]
Отсюда
\[
 s_0=\frac{u-F_0(f(u)-tx,g(u)-ty,h(u)-tz)}{t}
\]
является полиномом и $D(s_0)=1$. В компактной форме
\begin{align*}
 s_0={}&(-3f(u)^2-5g(u)-7)x+(h(u)-5f(u))y+(g(u)+2)z\\
       &+t\bigl(3f(u)x^2+5xy-yz\bigr)-t^2x^3.
\end{align*}
Эти тождества проверяются точной арифметикой в прилагаемом скрипте.

Однако $F_0$ не является координатой аффинного трёхмерного пространства:
\[
 \nabla F_0=(-3x^2-5y-7,\ z-5x,\ y+2)
\]
обращается в нуль в точках $(1,-2,5)$ и $(-1,-2,-5)$.
Координата полиномиального автоморфизма не имеет критических точек,
поскольку соответствующая строка обратимой матрицы Якоби не может
быть нулевой. Таким образом, полиномиальное восстановление параметра
и наличие слайса сами по себе не заменяют построения $F$.

Для нормированных полиномов Чебышёва $C_n(u)=2T_n(u/2)$ имеем
\[
 (C_3,C_4,C_5)=(f,g+2,h-5f).
\]
В этих аффинно изменённых координатах полином восстановления принимает
особенно простой вид
\[
 P_0(x,y,z)=yz-x^3+3x,\qquad P_0(C_3(u),C_4(u),C_5(u))=u.
\]
Это также непосредственно следует из
$C_4C_5=C_9+C_1$ и $C_3(C_3(u))=C_9(u)$.
Полином $P_0$ имеет критические точки $(1,0,0)$ и $(-1,0,0)$;
следовательно, даже столь короткая формула восстановления параметра
не является первой координатой выпрямляющего автоморфизма.

\subsection{Ещё два конкретных следствия и линеаризация действия}

Пусть $F=(P,Q,R_3)$ --- удостоверенный выпрямитель трилистника над
$K=\mathbb Q(\sqrt{-3})$, а $b(t)=t$. Приведённые формулы дают
$B=\ker D=K[t,p,q,r]$ и $A=B[s]$. Они непосредственно устраняют также
два связанных кандидата из того же первоисточника: Conjectures~5.2 и~5.3,
с.~1320. Ниже не используются методы гомологической алгебры.

\paragraph{Гиперповерхность из Conjecture~5.2.}
Обозначим через $\sigma=s_0$ исходный слайс, построенный из $F_0$.
Поскольку $D(\sigma)=D(s)=1$, разность
\[
 h_0=\sigma-s=\frac{P(V)-F_0(V)}t
\]
принадлежит $B$. Поэтому замена $s\mapsto s+h_0$ является
треугольным автоморфизмом $B[s]$, и
\[
 A=K[t,p,q,r,\sigma],\qquad A/(\sigma)=K[t,p,q,r]\cong K^{[4]}.
\]
Таким образом, $\sigma$ --- координата пятимерного пространства.
Это совместимо с некоординатностью $F_0$ на трёхмерном пространстве:
речь идёт о разных полиномах в разных кольцах.

\paragraph{Расширенное кольцо Риса из Conjecture~5.3.}
Положим $R_0=K[x,y,z]=K[P,Q,R_3]$ и $I=I_\gamma=(Q,R_3)$.
Рассматриваемое кольцо, определённое как подкольцо $R_0[t,t^{-1}]$,
имеет явное описание
\[
 R_0[t,I/t]=K[t,P,Q/t,R_3/t]\cong K^{[4]}.
\]
Действительно, старые координаты $Q,R_3$ восстанавливаются умножением
новых на $t$, а алгебраическая независимость новых четырёх образующих
следует после обращения $t$. После расширения скаляров до $\mathbb C$
оба изоморфизма сохраняются, давая отрицание именно этих двух
конкретных гипотез, как и Conjecture~5.1.

\paragraph{Естественное действие мультипликативной группы.}
Зададим действие подстановками
\[
 t\longmapsto\lambda t,\qquad
 u\longmapsto u,\qquad X\longmapsto\lambda^{-1}X
 \quad(\lambda\in K^*).
\]
Вектор $V=\gamma(u)-tX$ инвариантен. Поэтому на ядре $B$ новые
координаты $(t,p,q,r)$ имеют веса $(1,0,-1,-1)$. Следовательно,
это конкретное действие $\mathbb G_m$ на $\operatorname{Spec}B$
линеаризуется явной заменой координат. Общая проблема линеаризации
действий на аффинных пространствах отсюда не решается.

\subsection{Дихотомия Макар--Лиманова для связанного кольца}

Есть элементарное следствие, пригодное для дальнейшего поиска.
Пусть $B=\ker D_\gamma$ для
\[
 D_\gamma=t\partial_u+\sum_{i=1}^n f_i'(u)\partial_{x_i}
 \quad\text{на }k[t,u,x_1,\ldots,x_n].
\]
Тогда
\begin{equation}\label{eq:ml-dichotomy}
 \operatorname{ML}(B)\in\{k,\ k[t]\}.
\end{equation}
Здесь $\operatorname{ML}(B)$ --- пересечение ядер всех локально
нильпотентных $k$-дифференцирований $B$.

Действительно, $\partial_{x_i}$ коммутируют с $D_\gamma$, а потому
ограничиваются до LND на $B$. Следовательно,
\[
 \operatorname{ML}(B)\subseteq
 \bigcap_i\ker(\partial_{x_i}|_B)
 =B\cap k[t,u]=k[t].
\]
Если $\operatorname{ML}(B)$ содержит непостоянный полином $a(t)$, то
для всякого LND $\delta$ на $B$ имеем
$0=\delta(a(t))=a'(t)\delta(t)$. Кольцо является областью целостности,
характеристика равна нулю, поэтому $\delta(t)=0$. Значит $k[t]$
содержится в $\operatorname{ML}(B)$, что доказывает дихотомию.

Она задаёт два конкретных исследовательских направления. Достаточно
найти один LND $\delta$ на $B$ с $\delta(t)\ne0$, чтобы получить
$\operatorname{ML}(B)=k$. Если удастся доказать, что все LND уничтожают
$t$, то $\operatorname{ML}(B)=k[t]$, кольцо $B$ не полиномиально, и
соответствующее вложение невыпрямляемо. Обратная импликация
\emph{тривиальный ML $\Rightarrow$ полиномиальное кольцо} здесь не
утверждается. При известном выпрямителе формулы выше дают нужный LND
$\partial_t$ в координатах $(t,p,q,r)$.

\subsection{Сопряжённые LND на исходном пространстве}

Для проверенного $F=(P,Q,R)$ перенос координатных дифференцирований
даёт три коммутирующих LND $\delta_P,\delta_Q,\delta_R$ на
$k[x,y,z]$. В частности,
\[
 \delta_P(P)=1,\quad \delta_P(Q)=\delta_P(R)=0,
 \qquad\delta_P|_{\gamma}=\frac{d}{du}.
\]
Если $d=\det JF\in k^*$, то
\[
 \delta_P(a)=d^{-1}\det\frac{\partial(Q,R,a)}{\partial(x,y,z)}.
\]
Идеал кривой имеет вид $I_\gamma=(Q,R)$. Дифференцирования
$\delta_P$, $Q\delta_R$ и $R\delta_Q$ сохраняют этот идеал, и
пересечение их ядер равно $k$. Поэтому обычный инвариант
Макар--Лиманова самой $\mathbb A^1$ или окружающей $\mathbb A^3$
не различает рассматриваемые вложения. Содержательны условия,
связывающие LND с идеалом вложенной кривой, либо инварианты
связанного кольца $B$.

\subsection{Ограничения интерпретации и источники}

Невыпрямляемость вещественного узла не противоречит существованию
комплексного автоморфизма. Из комплексной полиномиальности ядра
нельзя без отдельного аргумента выводить вещественную или рациональную
полиномиальность. Также нельзя обращать импликацию
\emph{выпрямляемость $\Rightarrow$ полиномиальность ядра}, игнорируя
дополнительные условия на градуировку.

У K.~Masuda, \emph{Equivariant embeddings and semi-invariant locally
nilpotent derivations}, Theorems~1.1--1.2, рассматривается именно
эквивариантное соответствие и полиномиальные координаты с заданными
весами; это полезно для дальнейшего исследования симметрий гармонических
моделей. Доступный авторский текст:
\url{https://sci-tech.ksc.kwansei.ac.jp/~kayo/cmn/papers/eembed.pdf}.

Препринт S.~Oda \url{https://arxiv.org/abs/1201.4198}, появляющийся
в поиске как прежнее отрицательное решение гипотезы E, отозван автором
9 мая 2020 года с указанием ошибочности. Его утверждение не используется.

Первичный текст van den Essen--van Rossum доступен по ссылке издателя
\url{https://www.ams.org/journals/proc/2002-130-05/S0002-9939-01-06212-8/S0002-9939-01-06212-8.pdf};
проверенная копия того же журнального PDF:
\url{https://scispace.com/pdf/triangular-derivations-related-to-problems-on-affine-n-space-3val4k6v2j.pdf}.
